\documentclass[11pt]{article}

\usepackage[utf8]{inputenc}
\usepackage[T1]{fontenc}
\usepackage{lmodern}
\usepackage{geometry}
\usepackage{amsmath,amssymb,amsfonts,amsthm,mathtools}
\usepackage{bm}
\usepackage{enumitem}
\usepackage{microtype}
\usepackage{hyperref}
\usepackage{titlesec}
\usepackage{booktabs}
\usepackage{array}
\usepackage{setspace}
\usepackage{mathrsfs}
\usepackage{physics}

\geometry{margin=1in}
\hypersetup{colorlinks=true, linkcolor=black, urlcolor=blue, citecolor=black}
\setlist[enumerate]{topsep=0.4em,itemsep=0.35em}
\setlist[itemize]{topsep=0.3em,itemsep=0.25em}
\titleformat{\section}{\large\bfseries}{\thesection.}{0.5em}{}
\titleformat{\subsection}{\normalsize\bfseries}{\thesubsection.}{0.5em}{}
\setstretch{1.06}

\newcommand{\Aether}{\AE{}ther}
\newcommand{\Aflow}{\AE{}ther-flow}
\newcommand{\TheoryName}{\Aflow{} Interpretation of Relativity}
\newcommand{\TheoryProgram}{\Aether{} / \Aflow{} framework}
\newcommand{\ObserverMetric}{g^{\mathrm{br}}}
\newcommand{\CandidateMetric}{g^{\mathrm{cand}}}
\newcommand{\CompletedMetric}{g^{\sharp}}
\newcommand{\RedefMetric}{\widehat g}
\newcommand{\BaseShift}{\Sigma^{\mathrm{IER}}}
\newcommand{\ResponseShift}{\Sigma^{\mathrm{resp}}}
\newcommand{\CompShift}{\Sigma^{\mathrm{cmp}}}
\newcommand{\ResidualCore}{\mathcal V_{\mu\nu}^{\mathrm{res}}}
\newcommand{\ResidualDec}{\mathcal D_{\mu\nu}^{\mathrm{res}}}
\newcommand{\ReducedEin}{\mathcal L_{\ObserverMetric}}
\newcommand{\GaugeBook}{\mathcal G_{\ObserverMetric}}
\newcommand{\CompMix}{\mathcal M_{\mu\nu}^{\mathrm{cmp}}}
\newcommand{\CompQuad}{\mathcal Q_{\mu\nu}^{\mathrm{cmp}}}
\newcommand{\DeltaTComp}{\Delta T_{\mu\nu}^{\mathrm{cmp}}}
\newcommand{\ObsBurden}{\mathfrak B_{\mu\nu}^{\mathrm{obs}}}
\newcommand{\ObserverClass}[1]{\left[#1\right]_{\mathrm{obs}}}
\newcommand{\GaugeExact}{\mathfrak G_{\mu\nu}^{\mathrm{ex}}}
\newcommand{\ConstraintExact}{\mathfrak C_{\mu\nu}^{\mathrm{ex}}}

\newtheorem{definition}{Definition}
\newtheorem{proposition}{Proposition}
\newtheorem{remark}{Remark}
\newtheorem{corollary}{Corollary}
\newtheorem{theorem}{Theorem}

\title{\textbf{The \TheoryName{}}\\[0.4em]
\large Substrate-Response / Self-Response Charge-Polarization Candidate\\
\large Exact-Preservation Observer-Equation Exact Absorbability or Same-Package Field-Redefinition Attack Note}
\author{}
\date{}

\begin{document}

\maketitle

\begin{abstract}
The active one-metric exact-preservation line has already reached its sharpest currently proved observer-side theorem: the displayed burden
\[
\ObsBurden
=
\ResidualDec+\CompMix+\CompQuad-8\pi G_N\,\DeltaTComp
\]
does not survive as additional independent observer-side physics, so the completed branch is exactly observer-side effectively equivalent to Einstein--Hilbert plus ordinary matter through the same operative metric. That exact theorem still leaves one narrower bottleneck: can the remaining burden now be removed exactly, rather than merely shown non-independent?

The present note makes that bottleneck into the explicit next task. It fixes the ordered exact attack on \(\ObsBurden\). First, test direct exact absorbability: does \(\ObsBurden\) vanish in the exact observer quotient modulo the already-extracted gauge-exact and constraint-exact sectors? Second, if that direct route is not available on the recorded package, test an explicit same-package field redefinition: can the completed one-metric observer equation be rewritten exactly as Einstein--Hilbert plus ordinary matter without introducing a second operative metric or a second matter law?

The note proves the project-level success and failure tree attached to those tests. If either exact upgrade route succeeds, the branch acquires the strict observer-side closure theorem. If neither route succeeds on the fixed package, then the next honest paper is not another depth-for-its-own-sake continuation but a genuine necessity, maximality, or no-go theorem. A deeper-origin note is primary only if it directly supplies one of those missing exact upgrade ingredients. The present note does not yet prove either upgrade route; its positive role is to fix the exact next move and its honest outcomes.
\end{abstract}

\section{Introduction}

The observer-equation line has already passed through four exact tightenings on the same one-metric branch. The completed observer equation was written explicitly; the former core source was removed from that completed equation; the remaining displayed sectors were classified; and those sectors were then shown not to survive as additional independent observer-side physics. That sequence narrowed the exact burden sharply. It did not yet remove the displayed burden itself.

The project goal now forces a corresponding change of priority. The next manuscript should not pursue additional depth for its own sake while the live exact observer-equation bottleneck is already explicit. The next manuscript should attack that bottleneck directly.

\paragraph{Framework and claim boundary.}
\TheoryProgram{} denotes the broader \Aether{} / \Aflow{} framework built on the claims that the \Aether{} is the underlying four-dimensional substrate of reality, the \Aflow{} is its intrinsic ordered motion, observed three-dimensional space is the local experiential slice of that deeper substrate, S-time is the experienced order of change arising from matter, light, and the \Aflow{}, and observed expansion is the three-dimensional appearance of deeper four-dimensional ordered motion. Within that framework, \TheoryName{} denotes the exact-closure theory in which the effective gravitational dynamics are adopted to be exactly Einsteinian with universal matter coupling. In the active sequence, \emph{adoption} means use of that established relativistic dynamics without claiming substrate derivation, while \emph{derivation} is reserved for a first-principles recovery from explicit substrate variables.

The present note stays strictly inside that boundary. It does not claim a completed first-principles derivation of GR from substrate microphysics. It does not claim the strict identity \(\ObsBurden=0\). It does not claim that an exact same-package field-redefinition theorem has already been found. Its narrower role is to state the exact attack that now matters, and to fix the honest success and failure tree attached to that attack.

\section{Fixed Bottleneck and Strict Target}

\begin{definition}[Fixed one-metric exact-preservation package]
\label{def:fixed}
The present note keeps the following observer-level package fixed.
\begin{enumerate}
    \item The candidate and completed operative metrics are
    \begin{equation}
    \CandidateMetric
    =
    \ObserverMetric+\BaseShift+\ResponseShift,
    \qquad
    \CompletedMetric
    =
    \ObserverMetric+\BaseShift+\ResponseShift+\CompShift.
    \label{eq:metrics}
    \end{equation}
    \item The fixed completion solve is
    \begin{equation}
    \ReducedEin(\CompShift)=-\ResidualCore.
    \label{eq:solve}
    \end{equation}
    \item The completed observer equation is
    \begin{equation}
    G_{\mu\nu}[\CompletedMetric]
    =
    8\pi G_N\,T_{\mu\nu}^{\mathrm{matter}}[\CompletedMetric,\psi]
    +\GaugeBook(\CompShift)_{\mu\nu}
    +\ObsBurden,
    \label{eq:completed}
    \end{equation}
    with
    \begin{equation}
    \ObsBurden
    \equiv
    \ResidualDec+\CompMix+\CompQuad-8\pi G_N\,\DeltaTComp.
    \label{eq:burden}
    \end{equation}
    \item Gauge bookkeeping is already extracted explicitly in \eqref{eq:completed}, and the constraint-exact elimination sectors have already been removed before \(\ObsBurden\) is defined.
    \item The previously established exact observer-side theorem on this same branch is that the displayed sectors in \(\ObsBurden\) are exactly non-independent and therefore the completed branch is exactly observer-side effectively equivalent to Einstein--Hilbert plus ordinary matter through the same one operative metric.
\end{enumerate}
\end{definition}

\begin{definition}[Exact observer quotient]
\label{def:quot}
For the fixed package of Definition \ref{def:fixed}, the \emph{exact observer quotient} is the quotient of observer-side tensors by the already-extracted null classes:
\begin{enumerate}
    \item gauge-exact bookkeeping tensors;
    \item constraint-exact elimination tensors.
\end{enumerate}
The observer-quotient class of a tensor \(X_{\mu\nu}\) is denoted \(\ObserverClass{X_{\mu\nu}}\).
\end{definition}

\begin{remark}
Definition \ref{def:quot} does not enlarge the quotient by declaring exact non-independence itself to be zero. The present note keeps the null classes exactly as already extracted before the current attack begins.
\end{remark}

\begin{definition}[Strict observer-side closure on the fixed package]
\label{def:strict}
The fixed package is said to satisfy \emph{strict observer-side closure} if the completed observer equation admits an exact one-metric Einstein--matter presentation with no second operative metric and no second matter law, modulo only the already-extracted gauge-exact and constraint-exact sectors. Concretely, this means that either:
\begin{enumerate}
    \item the original completed metric \(\CompletedMetric\) itself satisfies
    \begin{equation}
    \ObserverClass{G_{\mu\nu}[\CompletedMetric]-8\pi G_N\,T_{\mu\nu}^{\mathrm{matter}}[\CompletedMetric,\psi]}=0,
    \label{eq:strictdirect}
    \end{equation}
    or
    \item there is an explicit exact invertible same-package rewriting through one rewritten metric variable \(\RedefMetric\) with the same single operative matter route, for which the rewritten observer equation has the same form in the quotient.
\end{enumerate}
\end{definition}

\begin{remark}
Strict observer-side closure is stronger than the already-proved effective-equivalence theorem. Effective equivalence says that the displayed burden does not survive as \emph{independent} observer-side physics. Strict closure asks whether the same burden can now be removed entirely, either directly or by an explicit same-package one-metric rewrite.
\end{remark}

\section{First Attack: Direct Exact Absorbability}

\begin{definition}[Direct exact absorbability]
\label{def:direct}
The burden \(\ObsBurden\) is \emph{directly exactly absorbable} on the fixed package if
\begin{equation}
\ObserverClass{\ObsBurden}=0.
\label{eq:obsvanish}
\end{equation}
Equivalently, there exist already-extracted gauge-exact and constraint-exact tensors \(\GaugeExact\) and \(\ConstraintExact\) such that
\begin{equation}
\ObsBurden=\GaugeExact+\ConstraintExact.
\label{eq:direct}
\end{equation}
\end{definition}

\begin{proposition}[Why direct absorbability is the first exact test]
\label{prop:first}
Direct exact absorbability is the first upgrade test on the fixed package because it asks whether the remaining burden already vanishes modulo the null sectors that have \emph{already} been extracted, without introducing any rewritten metric variable, second operative metric, or second matter law.
\end{proposition}

\begin{proof}
Definition \ref{def:direct} uses only the completed observer equation \eqref{eq:completed}, the displayed burden \eqref{eq:burden}, and the observer quotient of Definition \ref{def:quot}. It therefore stays entirely inside the recorded one-metric package and asks for no additional structure beyond an exact identity in that existing package. That makes it the minimal exact upgrade test still open on this line.
\end{proof}

\begin{theorem}[Direct exact absorbability gives strict observer-side closure]
\label{thm:direct}
If \(\ObsBurden\) is directly exactly absorbable in the sense of Definition \ref{def:direct}, then the fixed package satisfies strict observer-side closure in the sense of Definition \ref{def:strict}.
\end{theorem}

\begin{proof}
Insert \eqref{eq:direct} into the completed observer equation \eqref{eq:completed}. The terms beyond the Einstein--matter sector are then the already-extracted gauge-exact bookkeeping tensor together with gauge-exact and constraint-exact tensors. All of those vanish in the quotient of Definition \ref{def:quot}. Therefore \eqref{eq:strictdirect} holds, which is item (1) of Definition \ref{def:strict}.
\end{proof}

\section{Second Attack: Explicit Same-Package Field Redefinition}

\begin{definition}[Explicit same-package field-redefinition upgrade]
\label{def:redef}
An \emph{explicit same-package field-redefinition upgrade} on the fixed package means an explicit theorem exhibiting an exact invertible rewriting of the completed observer equation in terms of a rewritten metric variable \(\RedefMetric_{\mu\nu}\) such that:
\begin{enumerate}
    \item \(\RedefMetric_{\mu\nu}\) is determined only by the already fixed one-metric branch data;
    \item no second operative metric and no second matter law are introduced;
    \item the rewritten observer equation is exactly Einstein--Hilbert plus ordinary matter through that same rewritten one-metric route, modulo only the already-extracted gauge-exact and constraint-exact sectors.
\end{enumerate}
\end{definition}

\begin{remark}
The field-redefinition route is not permission to hide the burden inside a second operational law. The redefinition must stay in the same package: one operative metric, one universal matter route, and one exact observer-side Einstein--matter presentation after the rewrite.
\end{remark}

\begin{theorem}[An explicit same-package field redefinition gives strict observer-side closure]
\label{thm:redef}
If an explicit same-package field-redefinition upgrade in the sense of Definition \ref{def:redef} is established on the fixed package, then the fixed package satisfies strict observer-side closure in the sense of Definition \ref{def:strict}.
\end{theorem}

\begin{proof}
This is immediate from Definition \ref{def:redef}. The assumed rewritten metric variable \(\RedefMetric\) gives an exact one-metric Einstein--matter presentation with no second operative metric and no second matter law, modulo only the already-extracted null sectors. That is exactly item (2) of Definition \ref{def:strict}.
\end{proof}

\section{Ordered Next-Step Verdict}

\begin{theorem}[Ordered exact attack on \texorpdfstring{\(\ObsBurden\)}{ObsBurden}]
\label{thm:attack}
On the current one-metric exact-preservation branch, the immediate next exact move is the following ordered attack:
\begin{enumerate}
    \item test direct exact absorbability of \(\ObsBurden\) in the sense of Definition \ref{def:direct};
    \item if that direct route is not established on the fixed package, test an explicit same-package field-redefinition upgrade in the sense of Definition \ref{def:redef}.
\end{enumerate}
\end{theorem}

\begin{proof}
Proposition \ref{prop:first} identifies direct absorbability as the minimal exact test remaining on the fixed package. If that identity is not proved, Theorem \ref{thm:redef} identifies the remaining same-package route that can still upgrade the branch without leaving the one-metric observer-equation problem. Those are therefore the two ordered exact tests now standing between the current effective-equivalence theorem and strict observer-side closure.
\end{proof}

\begin{corollary}[Success route to the strict observer-side closure theorem]
\label{cor:success}
If either step of Theorem \ref{thm:attack} succeeds, then the branch acquires the strict observer-side closure theorem.
\end{corollary}

\begin{proof}
Step (1) implies the claim by Theorem \ref{thm:direct}. Step (2) implies the claim by Theorem \ref{thm:redef}.
\end{proof}

\begin{theorem}[Failure route to the next honest paper]
\label{thm:failure}
If neither direct exact absorbability nor an explicit same-package field-redefinition upgrade can be established on the fixed package, then the next honest primary paper is a genuine necessity, maximality, or no-go theorem for that same observer-side bottleneck.
\end{theorem}

\begin{proof}
By Corollary \ref{cor:success}, those two routes are exactly the current upgrade routes from effective equivalence to strict observer-side closure on the fixed package. If neither route can be established, then another manuscript that merely restates the same fixed package does not resolve the live bottleneck. The honest next primary task must then change from ``prove closure'' to ``prove why this package cannot be upgraded further,'' which is precisely the role of a genuine necessity, maximality, or no-go theorem.
\end{proof}

\begin{proposition}[Primary-status rule for deeper-origin continuation]
\label{prop:depth}
A deeper-origin note is primary on the present branch only if it directly supplies one of the missing exact upgrade ingredients identified in Theorem \ref{thm:attack}. A deeper-origin note that merely preserves the same completed observer equation without providing direct absorbability or an explicit same-package field redefinition is supporting context rather than the immediate next move.
\end{proposition}

\begin{proof}
The current bottleneck is already observer-equation exact and already written in the fixed package of Definition \ref{def:fixed}. Therefore a deeper-origin continuation resolves the live burden only if it changes that exact observer-side status by delivering one of the missing upgrade routes. If it does not, then it may still deepen the substrate story, but it does not discharge the exact bottleneck currently standing between effective equivalence and strict closure.
\end{proof}

\section{Conclusion}

The live exact observer-equation burden on the current one-metric exact-preservation branch is now fully ordered. The completed observer equation has already been reduced to the single displayed burden \(\ObsBurden\), and that burden has already been shown not to survive as additional independent observer-side physics. The next question is narrower: can \(\ObsBurden\) now be removed exactly?

This note fixes the correct attack. First test direct absorbability: whether \(\ObsBurden\) vanishes in the exact observer quotient modulo the already-extracted gauge-exact and constraint-exact sectors. If that route is not available on the recorded package, then test an explicit same-package field redefinition of the completed one-metric observer equation into exact Einstein--Hilbert plus ordinary matter without introducing a second operative metric or a second matter law.

If either route succeeds, the branch reaches the strict observer-side closure theorem. If neither route succeeds, then the next honest paper is a genuine necessity, maximality, or no-go theorem rather than another depth-for-its-own-sake continuation. A deeper-origin note is primary only if it directly contributes one of those missing exact upgrade ingredients. That is the exact next move now required by the project goal.

\end{document}
